%% LyX 2.4.4 created this file.  For more info, see https://www.lyx.org/.
%% Do not edit unless you really know what you are doing.
\documentclass[oneside,english]{amsart}
\usepackage[T1]{fontenc}
\usepackage[latin9]{inputenc}
\synctex=-1
\usepackage{amstext}
\usepackage{amsthm}

\makeatletter
%%%%%%%%%%%%%%%%%%%%%%%%%%%%%% Textclass specific LaTeX commands.
\newlength{\lyxlabelwidth}      % auxiliary length 

%%%%%%%%%%%%%%%%%%%%%%%%%%%%%% User specified LaTeX commands.
\date{}
\usepackage{commath}

\makeatother

\usepackage{babel}
\begin{document}

\section{Binary Lattices}

\label{section: binary lattices}

For positive integers $m$ and $n$, an $\left(m,n\right)$\emph{
lattice} is a rectangular grid that has $m+1$ rows and $n+1$ columns
of integers, which we call \emph{vertices}. All our lattices will
be \emph{binary}, meaning that each vertex is $0$ or $1$. A \emph{cell}
is a $\left(1,1\right)$ binary lattice, so that an $\left(m,n\right)$
binary lattice is made up of $mn$ cells. The sixteen possible cells
are 
\[
\begin{smallmatrix}0 & 0\\
0 & 0
\end{smallmatrix},\begin{smallmatrix}0 & 0\\
0 & 1
\end{smallmatrix},\begin{smallmatrix}0 & 0\\
1 & 0
\end{smallmatrix},\begin{smallmatrix}0 & 0\\
1 & 1
\end{smallmatrix},\begin{smallmatrix}0 & 1\\
0 & 0
\end{smallmatrix},\begin{smallmatrix}0 & 1\\
0 & 1
\end{smallmatrix},\begin{smallmatrix}0 & 1\\
1 & 0
\end{smallmatrix},\begin{smallmatrix}0 & 1\\
1 & 1
\end{smallmatrix},\begin{smallmatrix}1 & 0\\
0 & 0
\end{smallmatrix},\begin{smallmatrix}1 & 0\\
0 & 1
\end{smallmatrix},\begin{smallmatrix}1 & 0\\
1 & 0
\end{smallmatrix},\begin{smallmatrix}1 & 0\\
1 & 1
\end{smallmatrix},\begin{smallmatrix}1 & 1\\
0 & 0
\end{smallmatrix},\begin{smallmatrix}1 & 1\\
0 & 1
\end{smallmatrix},\begin{smallmatrix}1 & 1\\
1 & 0
\end{smallmatrix},\text{ and }\begin{smallmatrix}1 & 1\\
1 & 1
\end{smallmatrix}.
\]
 We further define that a lattice is \emph{framed} if it is binary
and every vertex on the boundary of the lattice is $0$, and a lattice
is \emph{proper} if it is framed and does not contain any $\begin{smallmatrix}0 & 1\\
1 & 0
\end{smallmatrix}$ or $\begin{smallmatrix}1 & 0\\
0 & 1
\end{smallmatrix}$ cells. For example, the right side of Figure~\ref{fig:example of f mapping}
is a $\left(5,4\right)$ proper lattice.

\end{document}
